% ---
% id: protocol-compiler/rem:coarse
% kind: rem
% title: "Direction of the outcome correspondence"
% summary: >-
%   With a trace-preserving chain, a correspondence from the emitted alphabet
%   onto the physical one forces every physical element covered more than once
%   to vanish, so it admits no coarse-graining. The correspondence therefore
%   runs from the physical alphabet onto the emitted one.
% prominence: standard
% status: accepted
% uses: [def:chaincond, thm:chain-consumption]
% assumes: []
% refs: []
% provenance: newly-derived
% lean:
% ---

\begin{remark}
\label{rem:coarse}
The direction in (C2) is forced, and the argument that forces it is a
\emph{typing} argument: (C2) fixes the direction of a tabulated map, it is
combinatorial, and it is read exactly.

Suppose instead a total surjection
$\gamma':\Vset\to\Vph$ from the emitted alphabet onto the physical one, with the
pullback read as $\Lambda^{\dagger}(F_v)=\Fph_{\gamma'(v)}$. With
trace-preserving channel data, Theorem~\ref{thm:chain-consumption}\,(iii)
says that this reading forces $\Fph_w=0$ for every outcome $w$ covered more
than once; the computation is Step~3 of its proof. The only non-injective $\gamma'$ that
survives is one whose extra outcomes are null, i.e.\ no coarse-graining at all.
Coarse-graining is a real operation --- a squashed family is typically
\emph{coarser} than the physical one --- so the correspondence must run from
the finer alphabet to the coarser, which is (C2), and the identity it feeds
must sum the fibre, which is \eqref{eq:pullback}. Emitting the pair
$(\gamma,\eqref{eq:pullback})$ rather than a bare correspondence of indices is
what keeps the statistics identity \eqref{eq:pullbackstats} true outcome by
outcome.
\end{remark}
